\documentclass[11pt]{article}
\usepackage[margin=1in]{geometry}
\usepackage{amsmath,amssymb,amsthm,booktabs,hyperref,microtype}
\newtheorem{proposition}{Proposition}
\newcommand{\Prb}{\mathbb P}
\title{Submission Audit: Parameter Scope and Sampling-Model Clarifications}
\author{}\date{October 22, 2026}
\begin{document}\maketitle
\noindent These notes supplement the earlier review of the refined
$(C,L_0)$ theorem. They address the original manuscript's optimal-order
statement and the integration of new constants into a standard sample
complexity theorem. They do not contradict the uniform sufficient bound
for the refined algorithm.

\section{A missing qualification on the matching lower bound}
The original manuscript states optimal complexity
$\Theta(n/(\epsilon\log n))$ without an explicit restriction on $\epsilon$
in Theorem 1.1. If interpreted as a two-parameter minimax statement with
universal constants on the whole interval $(0,1)$, this cannot be correct.
The issue is the matching lower bound, not the sufficient upper bound.

\begin{proposition}[Collision test near separation one]
For every integer $n\ge16$, let $\epsilon_n=1-n^{-2}$. There exists a
tester of $S\le n$ versus $\Delta_n\ge\epsilon_n$ using
$t=\lceil2\sqrt n\rceil+1$ independent samples and having both errors
less than $1/3$.
\end{proposition}
\begin{proof}
Return Yes when at least two samples coincide, and No otherwise.
Under the null, pad the distribution by zeros to a vector of length $n$.
The probability of all samples being distinct is $t!e_t(p_1,\ldots,p_n)$,
where $e_t$ is the elementary symmetric polynomial. It is maximized at
the uniform vector. One elementary proof chooses, among maximizers on
the simplex, a vector minimizing $\sum_i p_i^2$. For a pair of coordinates
$a,b$, with all others fixed,
\[
e_t=e_t(\mathrm{rest})+(a+b)e_{t-1}(\mathrm{rest})
                     +ab\,e_{t-2}(\mathrm{rest}).
\]
Replacing unequal $a,b$ by their average cannot decrease $e_t$ and strictly
decreases the sum of squares; hence the chosen maximizer must be uniform.
Therefore
\[
\Prb(\text{false No})\le \frac{(n)_t}{n^t}
\le \exp\left(-\frac{t(t-1)}{2n}\right)\le e^{-2}<\frac13.
\]
Here $t\le n$ for $n\ge16$, and the exponential inequality uses
$1-u\le e^{-u}$.

Under the alternative, let $q=\sum_{i=1}^n p_{(i)}\le n^{-2}$.
The tail probabilities satisfy $p_{(i)}\le q/n$ for $i>n$. Thus
\[
\sum_i p_i^2
\le q^2+\frac qn(1-q)
\le n^{-4}+n^{-3}.
\]
A union bound over sample pairs gives
\[
\Prb(\text{false Yes})\le {t\choose2}(n^{-4}+n^{-3})
\le\frac{25}{8}(n^{-3}+n^{-2})<\frac13,
\]
where $t\le2\sqrt n+2\le(5/2)\sqrt n$ for $n\ge16$.
The argument also holds for countable alphabets.
\end{proof}

Since $\epsilon_n\to1$, the asserted expression
$n/(\epsilon_n\log n)$ is of order $n/\log n$, while the tester above
uses $O(\sqrt n)=o(n/\log n)$ samples. Thus a universal matching lower
bound on all $(n,\epsilon)$ is ruled out. If the intended statement fixes
$\epsilon$ or bounds it above by a sufficiently small constant, that
restriction must be stated. Replacing $\log n$ by $\log(en)$ also removes
the notation's singularity at $n=1$ without changing asymptotic order.

\paragraph{A concrete route to a repaired statement.}
Retain the sufficient upper bound for all $n\ge1$, $0<\epsilon<1$.
State the matching lower bound separately with its justified parameter
range. For example, the already cited concurrent paper of Bao, Gao, and
Yao states in Theorem 4.4 a support-testing lower bound
$\Omega(s/(\delta\log s))$ for all sufficiently large $s$ and
$0<\delta<1/32$, on an alphabet of size $2s$, with fixed success above
$1/2$; Appendix A supplies the moment-matching and dilution argument.
See \href{https://arxiv.org/html/2609.37926v1#S4.SS1}{the source theorem}.
That result, if adopted with its hypotheses and attribution, is sufficient
for a matching lower-bound statement on $0<\epsilon<1/32$.
A broader parameter range should be proved, not inferred from an
unqualified citation to an estimation lower bound.

Wu--Yang Proposition 1 concerns support \emph{estimation} with a minimum
positive mass condition and restrictions on its own approximation
parameter. Converting it into the present support \emph{testing} problem
requires a separated pair of priors and the dilution step producing
$1/\epsilon$. Parameter symbols in the two problems must not be identified
without that reduction. See
\href{https://arxiv.org/pdf/1504.01227}{Wu--Yang, Proposition 1 and its model}.
Similarly, the original pretty-good-bound optimality reduction uses
$2\epsilon$ as a testing separation, so its stated range must be compatible
with a nonempty alternative and with the chosen lower-bound theorem.

\section{Expected samples versus a deterministic cap}
The refined theorem uses $M\sim\operatorname{Poi}(m)$ with
$m=Cn/(\epsilon\log(en))$. The coefficient 70 describes its mean and
does not impose a maximum sample count. The following coupling supplies
a standard capped-sample version without conditioning on completion.

\begin{proposition}[Cap conversion]
Let $b(C,L_0)$ be the computed certified common error endpoint in the new
certificate. Draw $M$ before querying the oracle. If
$M>K=\lceil2m\rceil$, return any allowed fixed outputs without queries;
otherwise run the original statistic. The capped algorithm uses at most
$K$ samples, expected sample use at most $m$, and joint error at most
\[
b(C,L_0)+\exp\{-C(2\log2-1)\}.
\]
For all eight certified parameter pairs this is still strictly below
their advertised rational error target.
\end{proposition}
\begin{proof}
Couple both algorithms using the same potential full sample sequence.
Their outputs agree unless $M>K$. Exponential Markov with parameter
$\log2$ gives $\Prb(M>K)\le e^{-m(2\log2-1)}$.
Because $n/\log(en)\ge1$ and $\epsilon<1$, $m>C$.
The error event for either required output in the capped experiment is
contained in the original joint error event union the overflow event.
No conditional independence claim is made. Finally, compare the displayed
sum with each target by outward Arb arithmetic; the results are saved in
\texttt{capped\_certificate.json}.
\end{proof}

Consequently $(70,14)$ can be stated as expected sample use at most
$70n/(\epsilon\log(en))$ and a deterministic maximum of
$\lceil140n/(\epsilon\log(en))\rceil$, with failure below $0.0009$.
It must not be described as a maximum-sample coefficient of 70.
Nor does the existing integer-minimality theorem prove a smallest possible
deterministic-cap constant over all algorithms.

\section{The scope of the refined necessary constant is different}
The new $C_{\rm crit}$ obstruction concerns one fixed Chebyshev tester,
with fixed finite $L_0>1$, over all input sizes and separations. Its
counterexample uses $\epsilon_n\to0$, not $\epsilon_n\to1$. It therefore
still excludes integer $C\le69$ even if the admissible separations are
restricted to $(0,1/32)$. There is no conflict with the collision example:
one is a limitation of a fixed construction in a small-separation regime;
the other concerns the global minimax rate in a large-separation regime.

The interval $(69.00627138,69.00627139)$ is a certified root bracket;
it does not prove an error lower bound for $C=69$ at the finite sizes
tested experimentally, or for an arbitrary fixed separation bounded away
from zero. The sufficient real constant is $144(\log2)^2$, so the exact
real optimum and the smallest sufficient $L_0$ remain open.

\section{Proof packaging before submission}
Changing the constants in the original algorithm block is not by itself
a valid integration. The original Lemma 4.1 uses $k_j<L/64$ to obtain
$\|f_j\|_\infty\le4n^{1/16}$. At $C=70$, scale zero instead has
$k_0/L\to\sqrt{512/70}\,3\log2>5$, so that premise fails.
The new coefficient/range envelopes and concentration proof must replace
the corresponding old argument, or the old theorem and the new
finite-constant theorem must be clearly stated as separate results.

The new concentration manuscript relies on the earlier expectation,
integral, coefficient, and half-line lemmas. A submission must include
these arguments or stable theorem-level references within the submitted
supplement; a relative link into a development directory is insufficient.
The infinite-support limit, Lindeberg verification, scale/degree coverage,
and scalar-root certificate should be separate named steps. Distinguish
analytic reasoning from finite arithmetic obligations: program-check counts
do not establish the mathematical reduction on their own.

Attribution should also preserve the distinction between the original
optimal-order theorem, concurrent results, refined constants for the same
statistic, and optional engineered hybrids. Neither a smaller constant
nor extensive simulation alone establishes a new general minimax lower
bound or guarantees acceptance at a particular conference.

\section{Why the current finite-size experiments cannot show binary gains}
For $n\le10^{12}$, $L=1+\log n<30$ and $70/L\ge2.4449$.
At mean $m=70n/(\epsilon L)$, the pure observed-count decision is already
uniformly very reliable. Its null error is zero; under $\Delta_r\ge\epsilon$
the discovery coupling gives
\[
\Prb(D\le r)\le\Prb\{\operatorname{Poi}(70n/L)\le n\}.
\]
On each quarter-width slab $L\in[a,b]\subset[1,30]$, exponential Markov
and $x-1-\log x\ge(x-1)^2/(2x)$ give the bound
\[
\exp\left\{-e^{a-1}\frac{(70/b-1)^2}{2(70/b)}\right\}.
\]
Outward rational arithmetic on all 116 slabs certifies their maximum to
be less than $1.863\times10^{-12}$. The exact endpoint is saved in
\texttt{finite\_counting\_certificate.json}. The same argument gives
the pretty-good coverage guarantee for the clipped count, by taking
$r=\min(e_\epsilon,n)-1$ when positive. It also gives simultaneous binary
and coverage correctness on the promises.

This explains the observed absence of any additional correct rejection
from the polynomial in the current $C=70$ experiments. At their maximum
sample size, the simple count already succeeds with an exceedingly small
uniform failure bound. More Monte Carlo repetitions or more distributions
in this same range will not yield a useful binary-power separation at the
same certified budget. The observed increase in the \emph{size} of the
support lower bound is a distinct and valid experimental outcome.
Reduced-budget exploration must be labelled separately from the uniformly
certified $C=70$ theorem. Conversely this finite-range count bound cannot
be extended to all $n$; it does not contradict the asymptotic need for
extrapolation or the necessary constant for the unchanged tester.
\end{document}
